\section{Setting, and how to read our claims}
\label{sec:setting}

\subsection{Finite systems with unbounded histories}

Every system we study has a finite physical state space, and finite sets of
actions, disturbances and outcomes. Runs, histories and event registers may be
unbounded, and so may the observations that report them. The temporal fields
include counters of order, progress and records, which grow without bound. Where they are,
our results about them are proved by invariants and induction, not by listing
cases. When a result concerns moving structure from one system to another,
the source and the target are distinct systems, fixed before anything is
evaluated. Re-encoding a single system does not count, and we show in
Section~\ref{sec:copy} why the distinction matters.

\subsection{The eight conditions}

We take the conditions from the closure architecture of
\citet{TsiokosArchitecture2026} and recall them only briefly. A system meets the
architecture when it satisfies the following eight conditions jointly, on one
system, with data tying them together:
\begin{enumerate}
\item \emph{objecthood}: staged closure kernels have robust fixed points
\item \emph{strict refinement}: some finer description is not recoverable from
the public one
\item a \emph{predictive self-coordinate}: refining the public description by
the system's own dynamics stabilizes at a predictive quotient
\item \emph{non-exposure}: the predictive state is not exposed at the given
public channels
\item \emph{action credit}: directionality is credited to self, other or none,
by where it is generated
\item \emph{temporal closure}: a funded order, tick and arrow carried by the
system's own record
\item \emph{binding}: joint content over branches strictly refines the
branch-by-branch description
\item \emph{viability}: a nonempty viability kernel under maintenance, which
underfunding empties.
\end{enumerate}

Each of our models meets these conditions at a recorded grade.
\begin{itemize}
\item A condition holds at \emph{design grade} when it is computed in the form
the architecture specifies, at an interface fixed in advance.
\item It holds in \emph{simplified form} when a declared stand-in replaces it.
\end{itemize}
On the two models that carry our strict continuation results, conditions (4)
to (7) hold only in simplified form:
\begin{itemize}
\item non-exposure is replaced by a single collision between two states that
look the same and predict differently
\item action credit by a declared credit rule
\item temporal closure by a declared order on events with a record and tick
added for the purpose
\item binding by a declared coupling.
\end{itemize}
The first replacement shows why these forms are weaker. A single collision
shows that the public channels cannot recover the whole predictive state.
Full non-exposure asks for more: that no public channel reveals any feature of
the predictive state that varies. On the two-copy model it fails, because the
public position is itself such a feature. The collision survives that
exposure, and full non-exposure does not.

On the third model, all eight conditions, as adopted at interfaces fixed in
advance, are computed jointly from a single execution, with stated choices for
conditions (1), (4), (5) and (6) (Section~\ref{sec:ancestry}). That is not a
certificate for the full architecture, whose complete panel of checks we do not
verify. Table~\ref{tab:grades} gives the details. We use the shorter phrase ``the target realizes the conditions'' only
with the grade stated nearby.

Realization in our sense is a check on the target's own states and
transitions. It resembles the functionalist sense in which a physical system
realizes a functional organization when its components and their dependencies
reflect those of the specification \citep{Chalmers1995Qualia}. The observation
that one organization can be realized by different systems goes back to
\citet{Putnam1967}. Whether passing the check is being conscious is not part
of the check.

% W14: source_reconciliation.md section 4 and ledger condition_grades.
\begin{table}[tbp]
\centering
\caption{The grade at which each of the eight conditions holds on each model. ``Computed'' means computed from the model's own dynamics. ``Simplified'' means a declared stand-in, described in Section~\ref{sec:setting}. ``Adopted interface'' means computed jointly from one execution at interfaces fixed in advance, with the noted choices.}
\label{tab:grades}
\small
\begin{tabularx}{\textwidth}{@{}>{\raggedright\arraybackslash}Xlll@{}}
\toprule
Condition & Two-copy model & Chain model & Fusion model \\
\midrule
(1) Objecthood & computed & computed & adopted interface\textsuperscript{a} \\
(2) Strict refinement & computed & computed & adopted interface \\
(3) Predictive self-coordinate & computed & computed & adopted interface \\
(4) Non-exposure & simplified & simplified & adopted interface\textsuperscript{b} \\
(5) Action credit & simplified & simplified & adopted interface\textsuperscript{c} \\
(6) Temporal closure & simplified & simplified & adopted interface\textsuperscript{c} \\
(7) Binding & simplified & simplified & adopted interface \\
(8) Viability & computed & computed & adopted interface \\
\bottomrule
\end{tabularx}
\par\smallskip
\begin{minipage}{\textwidth}
\footnotesize
\textsuperscript{a} Physical stationarity excludes growing histories; the period-two
two-step counterexample is retained.
\textsuperscript{b} Current observations on maintained support, not all future observations.
\textsuperscript{c} Completed local-cycle evidence with continuously authenticated
current residency and mandatory current-generator, ownership and control links.
\end{minipage}
\end{table}


\subsection{The temporal signature}

What a continuation must carry across, beyond the conditions themselves, is a
temporal organization. We fix it in advance as eight fields:
\begin{itemize}
\item the predictive present
\item an order on events
\item a tick phase
\item genuine progress
\item an arrow given by the growth of the system's own records
\item memory of the route by which the present was reached
\item a maintenance state
\item an owner.
\end{itemize}
We call these together the \emph{temporal signature}. It is fixed before, and
independently of, any judgement about lineage. The arrow in this list is an
ordering by record growth. It is not a thermodynamic arrow of time, and none
of our results concerns one.

\subsection{How to read our claims}

Our claims are of four kinds, and a result of one kind never supplies a claim
of another.
\begin{itemize}
\item \emph{Structural results}, theorems and their finite instances, are
proved in Lean.
\item \emph{Finite countermodels} refute a stated shortcut on an explicit
finite system, also checked in Lean. A countermodel refutes the claim it
represents on its stated boundary. It does not refute every way of
formulating the shortcut.
\item \emph{Interpretive hypotheses} link structure to experience, for example
that felt continuity requires the continuation of realized temporal closure.
We state them, mark them, and never use them as premises.
\item \emph{Ontological postulates}, such as the architecture's realization
postulate and Substrate Indifference, are premises the reader may grant or
decline. No result here uses them or supports them.
\end{itemize}

Words like transfer, upload, survival and same person appear only when we
describe others' thought experiments or say what our results do not show.
They never appear in the conclusion of a result.

Every theorem and countermodel is a declaration in the Lean development
that accompanies our results, built against Lean's core library with the standard
axioms only. Appendix~\ref{app:index} maps each result to its declarations.
The kernel checks the Lean statements. That those statements say what our
prose says is a separate matter. We checked it in a separate review, statement
by statement, described under Data and code availability.
